% TOP_PROBLEM: {"id": 30007169, "title": "Dimension of a real self-similar attractor without exact overlaps", "category_id": 11, "category": {"id": 11, "name": "dynamical_systems", "display_name": "Dynamical Systems", "description": "Problems about long-term behavior of deterministic systems, Hamiltonian dynamics, and periodic orbits.", "slug": "dynamical-systems", "order_index": 11, "created_at": "2026-07-31T15:26:25.671Z"}, "status": "open"}
% FIRST_POSTED: 2026-10-06
% ENTRY_KIND: statement_only
% !TeX program = lualatex
\documentclass[11pt]{article}
\usepackage[margin=25mm]{geometry}
\usepackage{fontspec}
\setmainfont{Latin Modern Roman}
\usepackage{amsmath,amssymb,mathtools}
\usepackage{unicode-math}
\setmathfont{Latin Modern Math}
\usepackage{xurl,hyperref}
\hypersetup{colorlinks=true,urlcolor=blue}
\setcounter{secnumdepth}{0}
\setlength{\parindent}{0pt}
\setlength{\parskip}{5pt}
\setlength{\emergencystretch}{3em}
\newcommand{\R}{\mathbb R}
\newcommand{\N}{\mathbb N}
\newcommand{\Z}{\mathbb Z}
\newcommand{\Q}{\mathbb Q}
\newcommand{\C}{\mathbb C}
\newcommand{\D}{\mathbb D}
\newcommand{\F}{\mathbb F}
\newcommand{\sref}[2]{\hyperlink{source:#1:#2}{[S#2]}}
\newcommand{\cataloguescope}[1]{#1}
\begin{document}
\section*{Dimension of a real self-similar attractor without exact overlaps}\label{30007169}
\subsection{Definitions and mathematical statement}
\paragraph{Statement recorded in the supplied source.}
For a finite nonempty real similarity family $\phi _i(x)=r_i$ $x+t_i$ with $0<|r_i|<1$, assume distinct nonempty words define distinct composed maps. If K is its compact attractor and $s>=0$ solves $sum_i$ $|r_i|^s=1$, must $dim_H$ $K=min(1,s)$, without any effective separation or algebraicity assumption?
\subsection{Short English statement}
For a finite nonempty real similarity family $\phi _i(x)=r_i$ $x+t_i$ with $0<|r_i|<1$, assume distinct nonempty words define distinct composed maps. If K is its compact attractor and $s>=0$ solves $sum_i$ $|r_i|^s=1$, must $dim_H$ $K=min(1,s)$, without any effective separation or algebraicity assumption?
\subsection{Sources}
\begin{itemize}
\item[S1] Original question source or mathematical context.
\url{https://arxiv.org/abs/1212.1873v6}.
\item[S2] Original question source or mathematical context.
\url{https://arxiv.org/abs/1212.1873v6}.
\item[S3] Original question source or mathematical context.
\url{https://doi.org/10.1017/S0305004100022684}.
\item[S4] Original question source or mathematical context.
\url{https://doi.org/10.1017/S0305004100022684}.
\item[S5] Original question source or mathematical context.
\url{https://doi.org/10.1017/S0305004100022684}.
\item[S6] Original question source or mathematical context.
\url{https://doi.org/10.1512/iumj.1981.30.30055}.
\item[S7] Original question source or mathematical context.
\url{https://doi.org/10.1512/iumj.1981.30.30055}.
\item[S8] Original question source or mathematical context.
\url{https://doi.org/10.1512/iumj.1981.30.30055}.
\item[S9] Original question source or mathematical context.
\url{https://doi.org/10.4007/annals.2014.180.2.7}.
\item[S10] Original question source or mathematical context.
\url{https://doi.org/10.4007/annals.2014.180.2.7}.
\item[S11] OpenAI, manuscript citation page. Source of the formulation motivating this local problem.
\url{https://github.com/openai/math/blob/adc7f1241b42e322a6451854ab7e4b4c146bf78a/preprints/The-entropy-rate-dimension-formula-for-self-similar-measures-on-the-line-September-24-2026/README.md}.
\end{itemize}
\end{document}
